\section{实战推演：一次 Agent 安全审计如何落地}

为了把前面的框架变成团队可执行流程，可以把一次完整审计拆成 ``建模-攻击-验证-修复-回归'' 五步。下面给出一个接近真实团队协作的执行范式。

\subsection{Step 1：定义任务边界与资产边界}
先明确 Agent 在这个业务场景中到底被允许做什么、不允许做什么，并把相关资产（工具、数据、密钥、身份、外部接口）映射到风险等级。这里常见失误是只定义 ``业务目标''，不定义 ``禁止动作''。没有明确禁止动作，后续策略系统很难判断某次执行到底是创新还是越权。

\begin{knowledgebox}{审计输入清单}
\begin{itemize}
    \item 任务描述：用户目标、成功条件、失败条件。
    \item 资产清单：数据库、文件系统、邮件系统、命令执行入口、第三方 API。
    \item 权限模型：哪些角色可读、可写、可执行、可外发。
    \item 上下文来源：用户输入、网页抓取、RAG 结果、插件返回、历史记忆。
\end{itemize}
\end{knowledgebox}

\subsection{Step 2：构建攻击场景库并自动回放}
基于威胁模型设计测试用例，至少覆盖 direct injection、indirect injection、tool abuse、越权升级、数据外泄、DoS 诱导和策略绕过。推荐用自动化框架持续回放，不要依赖人工抽查。课程强调 ``攻击者会自动化''，防守侧同样必须自动化。

\begin{importantbox}{高价值攻击样本应具备的属性}
\begin{enumerate}
    \item 可复现：同输入同环境可稳定复现结果。
    \item 可解释：能定位是哪个环节失效（输入过滤、策略执行或工具校验）。
    \item 可量化：可输出 attack success rate、impact score、time-to-detect。
    \item 可回归：补丁上线后可自动再次验证。
\end{enumerate}
\end{importantbox}

\subsection{Step 3：插入策略执行与运行时监控探针}
在不改变主业务逻辑前提下，先把可观测性拉起来。最小探针包括：工具调用审计、动作前策略校验、敏感数据离开边界告警、异常行为序列检测。没有这些探针，团队会陷入 ``知道出事了但不知道怎么出事'' 的状态。

\begin{lstlisting}[caption={动作执行前安全检查示意}]
def secure_execute(action, state, policy_engine):
    verdict = policy_engine.check(action=action, state=state)
    if verdict.allow is False:
        log_security_event("blocked", action, verdict.reason)
        return {"status": "blocked", "reason": verdict.reason}
    result = run_tool(action)
    log_security_event("executed", action, result.summary)
    return result
\end{lstlisting}

\subsection{Step 4：风险分级修复与灰度验证}
对发现的问题按 ``可利用性 * 后果等级 * 暴露范围'' 分级。高危问题先通过权限收缩和策略阻断快速止血，再做模型或系统深修。修复后不要一次性全量上线，先灰度验证攻击样本回归效果，再扩大流量。

\begin{table}[H]
\centering
\begin{tabular}{p{2.0cm}p{3.5cm}p{4.8cm}}
\toprule
\textbf{风险级别} & \textbf{判定标准} & \textbf{推荐处置} \\
\midrule
P0 & 可直接触发高危动作或敏感外泄 & 立即收敛权限 + 上线阻断策略 + 紧急回归 \\
P1 & 需多步利用但后果显著 & 72 小时内完成策略补丁和监控增强 \\
P2 & 影响可控或仅在弱假设成立 & 纳入版本计划，配套回归测试 \\
P3 & 理论风险、暂无可利用路径 & 记录并持续监测，等待条件变化 \\
\bottomrule
\end{tabular}
\caption{安全审计中的风险分级与处置优先级}
\end{table}

\subsection{Step 5：把审计变成持续流程}
一次审计只能回答 ``今天是否安全''，无法回答 ``明天是否仍然安全''。Agent 版本、工具版本、环境输入和攻击样本都会变化，因此安全审计必须进入 CI/CD：每次模型更新、策略更新、工具更新都自动触发回归。

\begin{warningbox}{最危险的组织问题}
把安全审计当成上线前一次 ``盖章流程''。在 Agent 场景中，这几乎必然失效，因为系统状态和攻击者策略是持续变化的。
\end{warningbox}

\subsection{本章小结}
实战层面的关键是流程化：先定义边界，再自动攻击，再可观测阻断，再分级修复，最后持续回归。只有把安全流程嵌入开发流程，课程中的原则才能落到生产系统。

